
% =============================================================================
% Section 01 — The Ontological Problem of Invariance
% Book I — Ontology of Viability
% Chapter 02 — Invariance
%
% Purpose:
%   To establish invariance as an independent ontological problem
%   and distinguish it from identity, persistence, continuity,
%   stability, and immutability.
%
% Conceptual Policy:
%   1. Invariance is not a definition or criterion of identity.
%   2. Invariance is not equivalent to the absence of change.
%   3. The existence of an invariant is not universally presupposed.
%   4. Every claim of invariance requires specification of its
%      domain and respect.
%   5. Ontological invariance is not reducible to epistemic
%      recognition, measurement, or description.
% =============================================================================

\section{The Ontological Problem of Invariance}
\label{sec:ontological-problem-of-invariance}

\subsection{The Problem of What Remains}

Existence is not necessarily characterized by an absence of change.
An existent may undergo transformations in its determinations,
configuration, composition, relations, or manifestations while
remaining an existent.

The occurrence of change, however, does not by itself establish
that everything associated with an existent changes. Nor does
the persistence of some determination or relation establish
that the existent as a whole remains unchanged.

This raises a distinct ontological problem:

\begin{quote}
What, if anything, remains invariant in an existent when it
undergoes transformation, and under what conditions can such
invariance be said to obtain?
\end{quote}

The problem concerns neither change alone nor the mere persistence
of an existent. It concerns the possible preservation of a
determination, relation, configuration, or other ontologically
relevant aspect across a specified transformation.

Invariance therefore requires an analysis of the relation between
what undergoes transformation and what, within the relevant
scope of that transformation, remains preserved.

The question is not whether an existent is simply the same
before and after a change. Such a formulation would risk
presupposing an answer to the problem of identity before
the relevant ontological conditions have been established.

Rather, the question concerns the possible preservation of
particular aspects of an existent under transformation,
without assuming that such preservation is equivalent to
the persistence of the existent's identity.

\subsection{Invariance as an Independent Ontological Problem}

Invariance is often discussed in connection with mathematical
structures, physical laws, symmetries, conservation principles,
and the behavior of systems under transformation.

These contexts provide important analytical resources.
They demonstrate that different properties, relations, or
structural features may remain preserved under specified
transformations.

However, the existence of such formal or scientific descriptions
does not, by itself, settle the ontological status of invariance.

A mathematical invariant is defined relative to a mathematical
structure and a specified class of transformations.
A physical invariant is established within a physical theory
and its relevant conditions.
A structural invariant may concern the preservation of an
organizational pattern despite changes in its components.

These examples indicate different possible forms of preservation,
but they do not establish that all such forms possess the same
ontological status or can be reduced to a single underlying
mechanism.

The ontological problem therefore requires a distinction
between the formal identification of an invariant and the
ontological interpretation of what is preserved.

An invariant cannot be regarded as ontologically established
merely because a quantity, relation, or expression remains
formally unchanged within a particular representation.

The relevant question is whether the preservation corresponds
to an actual determination or relation of the existent,
rather than exclusively to a feature of the formal framework
used to describe it.

This distinction does not diminish the significance of
mathematical or scientific results. It identifies the further
ontological question that must be addressed before those
results can support claims about the nature of existence.

\subsection{The Relational Character of Invariance}

Invariance cannot be adequately understood as an isolated
property attributed to an existent without reference to
the conditions under which preservation is asserted.

A claim of invariance necessarily involves several analytically
distinguishable elements:

\begin{enumerate}
    \item \textbf{The subject of invariance:}
    The existent, determination, relation, or configuration
    whose possible preservation is under consideration.

    \item \textbf{The transformation:}
    The change or transformation against which preservation
    is assessed.

    \item \textbf{The domain of transformation:}
    The range of conditions, states, configurations, or
    circumstances within which the transformation is considered.

    \item \textbf{The respect of invariance:}
    The particular determination, relation, magnitude,
    or organizational feature that is claimed to remain
    preserved.
\end{enumerate}

These elements are analytically distinguishable but need not
correspond to four independently existing entities.

In particular, the domain of transformation and the respect
of invariance specify the scope and terms of the claim.
They do not necessarily constitute additional ontological
components of the existent.

The relational character of invariance can be expressed
provisionally as follows:

\begin{quote}
An existent or one of its determinations is invariant
with respect to a specified transformation when the
relevant determination or relation remains preserved
throughout the transformation within the specified domain
and under the specified respect.
\end{quote}

This formulation is provisional because the ontological
meaning of preservation, the admissibility of transformations,
and the status of the relevant domain require further analysis
in the subsequent sections.

It also leaves open whether invariance should ultimately
be understood as a relation, a condition of preservation,
a feature of an existent, or a more complex ontological
configuration.

The relational formulation is therefore an analytical
starting point rather than a final ontological definition.

\subsection{Invariance and the Possibility of Change}

The existence of change does not necessarily exclude
invariance.

An existent may undergo changes in some determinations
while preserving others. Its composition may change while
a specified structural relation remains preserved.
Its quantitative magnitudes may vary while a particular
qualitative determination remains unchanged.
Its components may be replaced while a certain organizational
pattern continues to obtain.

These possibilities demonstrate that change and invariance
are not necessarily contradictory.

Their compatibility, however, depends on the respect under
consideration.

A determination may be invariant under one transformation
and non-invariant under another. A relation may be preserved
within one domain but cease to obtain beyond its boundaries.
A configuration may remain invariant under certain changes
while being altered by others.

Consequently, an assertion of invariance without specifying
the relevant transformation and domain is ontologically
incomplete.

The same applies to the opposite assertion. The fact that
an existent undergoes transformation does not justify
the conclusion that every determination, relation, or
structural feature associated with it has changed.

Invariance and change must therefore be examined in relation
to one another without reducing either to the other.

The relation between an initial configuration, a transformation,
and a resulting configuration can be represented schematically.
The representation does not presuppose that an invariant
must exist. It indicates that preservation and variation
are possibilities to be assessed with respect to specified
determinations and within a specified scope.


The figure illustrates that a transformation may be associated
with both possible preservation and possible variation.
Neither outcome should be assumed in advance. Whether a
determination remains invariant must be established through
an analysis of the relevant transformation, domain, and
respect of preservation.

\subsection{Distinguishing Invariance from Related Ontological Concepts}

The ontological problem of invariance requires distinguishing
it from several concepts that are frequently associated
with preservation and persistence.

These concepts may overlap in particular cases, but they
are not interchangeable.

\paragraph{Invariance and Identity.}

Identity concerns the ontological identity of an existent.
Invariance concerns the possible preservation of a specified
determination or relation under transformation.

The preservation of an invariant does not, by itself,
establish the numerical identity of an existent.
Likewise, the absence of an identified invariant does not,
by itself, establish the absence of identity.

Invariance therefore cannot serve as a general definition
or sufficient criterion of identity.

\paragraph{Invariance and Persistence.}

Persistence concerns the continued existence of an existent
through some interval or succession of conditions.

Invariance concerns what remains preserved under a specified
transformation.

An existent may persist while undergoing changes in many
of its determinations. Conversely, the preservation of
a particular invariant does not independently establish
the persistence of the existent as a whole.

\paragraph{Invariance and Continuity.}

Continuity concerns the manner in which existence, determination,
or transformation proceeds without a relevant rupture or
discontinuity.

Invariance concerns preservation under transformation.

A continuous process may involve extensive change, while
a discontinuous transformation may preserve certain
determinations or relations.

The two concepts therefore address different ontological
questions.

\paragraph{Invariance and Stability.}

Stability concerns the maintenance or boundedness of a
condition, configuration, or organization under relevant
conditions or disturbances.

Invariance concerns the preservation of a specified feature
under a specified transformation.

Stability may involve resistance to change, restoration,
or maintenance within a range of variation.
Invariance need not imply resistance, restoration, or
dynamic equilibrium.

\paragraph{Invariance and Immutability.}

Immutability denotes the absence of change in the relevant
respect.

Invariance does not necessarily require the absence of
change in the existent as a whole. It is compatible with
transformation, provided that the specified determination
or relation remains preserved.

Thus, invariance must not be reduced to immutability.

These distinctions establish that invariance occupies
a distinct analytical position within ontology.
Its relation to each of these concepts must be investigated
without assuming conceptual equivalence.

The distinction between invariance and its neighboring
ontological concepts does not imply that invariance itself
has only one form. Different determinations and relations
may be preserved under transformation, giving rise to
different analytical forms of invariance.

The following classification identifies several such forms.
It is an analytical classification rather than an assertion
that these forms are mutually exclusive, exhaustive, or
ontologically independent.


The forms presented in the table identify different respects
in which preservation may be examined. Their classification
does not establish that every existent exhibits all these
forms, nor that any particular form must obtain under every
transformation.

\subsection{The Scope and Limits of an Invariance Claim}

An assertion that something is invariant is incomplete
unless its scope is specified.

At minimum, such an assertion requires clarification of:

\begin{itemize}
    \item what is claimed to remain preserved;
    \item which transformation is relevant;
    \item which domain of transformation is considered;
    \item under which respect preservation is asserted;
    \item whether preservation is absolute within that domain
    or conditional upon particular circumstances.
\end{itemize}

Without these specifications, invariance risks becoming
an unrestricted attribution of permanence.

No universal assumption should be made that every existent
possesses an invariant, that every transformation preserves
some feature, or that every apparent preservation has
ontological significance.

An invariant may exist only under restricted conditions.
It may be preserved within one domain and cease to be
preserved beyond it.
Different determinations of the same existent may exhibit
different patterns of preservation.

Furthermore, the identification of an invariant within
a limited domain does not justify its generalization
to all possible transformations.

The scope of an invariance claim must therefore remain
proportionate to the conditions under which preservation
is established.

\subsection{The Ontological Problem and Its Subsequent Development}

The preceding analysis establishes invariance as a distinct
ontological problem concerning the possible preservation
of determinations and relations under transformation.

It does not yet establish a complete definition of invariance,
a final classification of its forms, or a universal account
of the conditions under which invariance obtains.

Those questions require separate examination.

The subsequent sections will develop the problem through
the following sequence:

\begin{enumerate}
    \item The relation between invariance and identity,
    with particular attention to their ontological distinction.

    \item The nature and defining conditions of invariants,
    including the criteria by which preservation may be
    distinguished from formal or apparent sameness.

    \item The relation between invariance and transformation,
    including the conditions under which change and
    preservation can coexist.

    \item The relation between invariance and difference,
    with attention to the possibility of preserved features
    within configurations that undergo differentiation.

    \item The role of structure and composition in the
    preservation of determinations and relations.

    \item The relation between invariance and continuity,
    without reducing either concept to the other.

    \item The forms and domains within which invariance
    may obtain.

    \item The limits of invariance and the conditions
    under which claims of preservation cease to be valid.
\end{enumerate}

The purpose of this development is not to establish
invariance as the ultimate foundation of identity or
existence.

Rather, it is to determine the ontological status,
scope, conditions, and limitations of preservation
within a world in which transformation is possible.

The central question remains:

\begin{quote}
What can remain invariant through transformation,
and what does such preservation mean for the ontological
constitution of the existent?
\end{quote}
